\documentclass[12pt,a4paper,oneside]{report}

%==============================================================================
% PACKAGES AND CONFIGURATION (ENSSEA STRICT COMPLIANCE)
%==============================================================================
\usepackage[utf8]{inputenc}
\usepackage[T1]{fontenc}
\usepackage[english]{babel}
\usepackage[a4paper, top=2cm, bottom=2cm, left=3cm, right=2cm]{geometry}
\usepackage{setspace}
\usepackage{fancyhdr}
\usepackage{graphicx}
\usepackage{float}
\usepackage{array}
\usepackage{booktabs}
\usepackage{caption}
\usepackage{subcaption}
\usepackage{xcolor}
\usepackage{hyperref}
\usepackage{titlesec}
\usepackage{amsmath}
\usepackage{amssymb}
\usepackage[bottom]{footmisc}
\usepackage{tabularx}
\usepackage{multirow}
\usepackage{longtable}
\usepackage{colortbl}

%==============================================================================
% TYPOGRAPHY
%==============================================================================
\usepackage[scaled=0.92]{helvet}
\renewcommand{\familydefault}{\sfdefault}

\setstretch{1.5}
\setlength{\parindent}{0pt}
\setlength{\parskip}{0.6cm}

%==============================================================================
% TITLE FORMATTING
%==============================================================================
\titleformat{\chapter}[display]
  {\centering\bfseries\fontsize{20pt}{24pt}\selectfont}
  {Chapitre \thechapter}
  {1em}
  {}

\titleformat{\section}[hang]
  {\bfseries\fontsize{14pt}{18pt}\selectfont}
  {\thesection}
  {1em}
  {}

\titleformat{\subsection}[hang]
  {\bfseries\underline\fontsize{12pt}{14pt}\selectfont}
  {\thesubsection}
  {1em}
  {}

\titleformat{\subsubsection}[hang]
  {\bfseries\fontsize{12pt}{14pt}\selectfont}
  {\thesubsubsection}
  {1em}
  {}

%==============================================================================
% PAGE STYLE
%==============================================================================
\pagestyle{fancy}
\fancyhf{}
\fancyfoot[C]{\thepage}
\renewcommand{\headrulewidth}{0pt}

\begin{document}

%==============================================================================
% CHAPTER I: THEORETICAL FOUNDATIONS
%==============================================================================
\chapter{Theoretical Foundations: Data-Driven Marketing \& Retail}

% ------------------------------------------------------------------------------
% INTRODUCTION TO CHAPTER I
% ------------------------------------------------------------------------------
\section*{Introduction to Chapter I}
\addcontentsline{toc}{section}{Introduction to Chapter I}

The contemporary retail landscape is experiencing unprecedented transformation driven by the exponential growth of data generation, advanced computational capabilities, and evolving consumer expectations. Organizations operating in the retail sector face mounting pressure to transition from traditional, intuition-based decision-making frameworks toward sophisticated, evidence-based strategies grounded in comprehensive customer data analysis. This paradigm shift, commonly referred to as data-driven marketing, represents one of the most significant transformations in modern business strategy and operational execution.

Data-driven marketing, as a strategic discipline, fundamentally reshapes how retailers understand, engage, and retain their customer base. By leveraging quantitative and qualitative data sources—ranging from transactional histories and behavioral patterns to demographic characteristics and psychographic profiles—retailers can develop increasingly granular and actionable insights into customer preferences, purchasing patterns, and lifetime value potential. The integration of these insights into strategic decision-making processes enables organizations to optimize resource allocation, enhance customer experience personalization, and ultimately drive sustainable competitive advantage.

The imperative for customer segmentation emerges naturally within this context. As organizations accumulate increasingly vast and complex datasets, the challenge of identifying meaningful patterns and extracting actionable intelligence becomes correspondingly more demanding. Customer segmentation—the process of dividing a heterogeneous customer population into homogeneous subgroups based on shared characteristics or behaviors—provides a structured analytical framework through which organizations can navigate this complexity and translate raw data into strategic value. This chapter, therefore, provides a comprehensive exploration of the theoretical foundations, methodological approaches, and practical applications of customer segmentation within data-driven retail contexts.

This introductory chapter is organized into three primary sections, each addressing distinct yet interconnected dimensions of data-driven marketing and customer segmentation. The first section examines the conceptual underpinnings, strategic relevance, and evolutionary trajectory of customer segmentation from traditional approaches toward contemporary data-driven methodologies. The second section critically evaluates traditional segmentation paradigms—including Recency-Frequency-Monetary (RFM) analysis, demographic segmentation, and behavioral segmentation—discussing their underlying assumptions, analytical procedures, and practical applications within retail contexts. The final section addresses the inherent limitations of classical segmentation methodologies when confronted with the complexity, volume, and velocity of contemporary retail data ecosystems, thereby establishing the conceptual foundation for subsequent chapters addressing machine learning-based approaches.

% ==============================================================================
% SECTION 1.1
% ==============================================================================
\section{1.1 — Client Segmentation: Concepts, Issues, and Evolution}

This section provides a comprehensive examination of customer segmentation as a strategic and analytical discipline, tracing its conceptual evolution from foundational principles through contemporary manifestations. The discussion encompasses the formal definition of customer segmentation, its strategic importance within retail and marketing contexts, and the historical development of segmentation approaches in response to changing market conditions and technological capabilities.

\subsection{1.1.1 — Definition and Strategic Importance of Customer Segmentation}

\subsubsection{Foundational Definition and Conceptual Framework}

Customer segmentation, in its most general form, refers to the process of partitioning a heterogeneous customer population into smaller, relatively homogeneous subgroups—commonly termed "segments"—based on specific criteria or characteristics deemed relevant to the organization's strategic objectives.\footnote{Kotler, P., Keller, K. L. (2016): \textit{Marketing Management}, 15th Edition, Pearson Education, New Jersey, p. 234.} This partitioning process fundamentally acknowledges that customer populations are intrinsically diverse, exhibiting significant variation across multiple dimensions including purchasing behavior, product preferences, price sensitivity, communication channel preferences, and lifetime value potential.\footnote{Wind, Y. (1978): \textit{Issues and Advances in Segmentation Research}, Journal of Marketing Research, Vol. 15, No. 3, pp. 317-337, p. 322.}

The analytical objective underlying segmentation endeavors is to identify and extract patterns of similarity within the larger customer population, thereby enabling organizations to develop differentiated strategies tailored to the specific characteristics, needs, and preferences of distinct customer segments. This differentiation approach stands in marked contrast to mass marketing paradigms that implicitly assume a relatively homogeneous customer population with relatively uniform preferences and needs.

The theoretical foundation for customer segmentation rests upon several fundamental premises: First, customer populations are heterogeneous with respect to relevant characteristics and behavioral patterns. Second, measurable variation exists across these dimensions, enabling empirical identification of meaningful groupings. Third, segments thus identified remain relatively stable over defined temporal periods, providing sufficient continuity for strategic planning and resource allocation. Fourth, segments exhibit differential responsiveness to marketing stimuli and strategies, enabling organizations to achieve superior outcomes through segment-specific strategy development and implementation.\footnote{Wedel, M., Kamakura, W. A. (2000): \textit{Market Segmentation: Conceptual and Methodological Foundations}, 2nd Edition, Kluwer Academic Publishers, Boston, p. 18.}

The practical importance of customer segmentation within retail and marketing contexts is multifaceted and substantial. Segmentation enables retailers to:
\begin{itemize}
    \item Develop targeted marketing communications that resonate with specific preferences and communication channel preferences.
    \item Optimize product assortment and pricing strategies to align with segment-specific demand patterns.
    \item Allocate limited resources—including marketing budgets and inventory—in a manner that reflects differential profitability.
    \item Identify and prioritize high-value customers who warrant intensive relationship development and retention efforts.
    \item Develop personalized customer experience strategies that enhance satisfaction, loyalty, and repurchase behavior.
\end{itemize}

\subsection{1.1.2 — Traditional Approaches to Customer Segmentation}

This subsection examines the principal traditional segmentation methodologies employed within marketing and retail contexts prior to the widespread adoption of machine learning techniques.

\subsubsection{1.1.2.1 — Demographic Segmentation}

Demographic segmentation partitions customer populations based on observable characteristics including age, gender, income level, education, family structure, and geographic location.\footnote{Schiffman, L. G., Kanuk, L. L. (2010): \textit{Consumer Behavior}, 10th Edition, Prentice Hall, New Jersey, pp. 45-67.} The widespread adoption of demographic segmentation stems from readily available data sources and the ease of understanding for stakeholders.

Despite these advantages, demographic segmentation exhibits significant limitations. Most critically, demographic characteristics exhibit weak correlation with consumer preferences and purchasing behavior in many retail categories.\footnote{DeCarlo, T. E. (2005): \textit{The Effects of Sales Message and Suspicion of Ulterior Motives on Salesperson Evaluation}, Journal of Consumer Psychology, Vol. 15, No. 3, pp. 238-249.} Individuals within identical demographic categories frequently exhibit substantial variation in brand loyalty and spending patterns.

\subsubsection{1.1.2.2 — Recency, Frequency, and Monetary (RFM) Analysis}

RFM analysis represents one of the most widely utilized traditional segmentation methodologies within retail. It partitions the customer population based on three behavioral dimensions extracted from transactional data:
\begin{itemize}
    \item \textbf{Recency (R)}: The temporal interval between the reference date and the customer's most recent transaction.
    \item \textbf{Frequency (F)}: The total number of transactions completed by a customer within a defined period.
    \item \textbf{Monetary (M)}: The total monetary value (revenue) generated by a customer.
\end{itemize}

The theoretical foundation rests upon the observation that recent, frequent, and high-value customers exhibit the highest probability of future purchase behavior.\footnote{Hughes, A. M. (2005): \textit{Strategic Database Marketing}, 3rd Edition, American Marketing Association, Chicago, p. 156.} RFM segmentation offers several advantages: it requires only basic transactional data, is straightforward to interpret, and correlates robustly with customer profitability.

However, RFM segmentation captures only previous transactional behavior and ignores important non-behavioral dimensions including customer attitudes and psychological characteristics. Furthermore, it assumes that historical behavior remains predictive of future behavior—an assumption that may not hold in volatile markets.

\subsubsection{1.1.2.3 — Behavioral Segmentation}

Behavioral segmentation extends beyond RFM by incorporating product category preferences, brand preferences, shopping channel preferences (online vs. physical), and purchase occasion characteristics.\footnote{Plummer, J. T. (1974): \textit{The Concept and Application of Life Style Segmentation}, Journal of Marketing, Vol. 38, No. 1, pp. 33-37.} It allows retailers to identify distinct segments such as "Online Channel Specialists" or "Promotional Hunters."

While behavioral characteristics typically exhibit stronger correlation with profitability than demographics, these approaches capture \textit{what} customers do but provide limited insight into \textit{why} they engage in specific behaviors.

\subsection{1.1.3 — Limitations of Classical Methods in Addressing Retail Data Complexity}

Contemporary retail environments are characterized by complexity across data volume, velocity, and behavioral heterogeneity that exceed the analytical capacity of traditional methodologies.

\subsubsection{1.1.3.1 — The Curse of Dimensionality}

Traditional methodologies rely on relatively small numbers of variables (3 to 10). Contemporary retail environments generate data across substantially higher dimensional spaces (transaction timing, browsing patterns, engagement metrics, contextual factors).\footnote{Bellman, R. E. (1961): \textit{Adaptive Control Processes: A Guided Tour}, Princeton University Press, p. 76.} Traditional clustering algorithms encounter computational challenges in such spaces, and distance metrics become increasingly problematic as relative distances converge—a phenomenon known as the "curse of dimensionality."

\subsubsection{1.1.3.2 — Behavioral Heterogeneity and Non-Stationarity}

Retail customers increasingly engage in "channel switching" and "cross-channel shopping," violating implicit assumptions of stability. Customer preferences exhibit substantial volatility in response to factor such as life stage transitions and competitive pressure.\footnote{Guadagni, P. M., Little, J. D. (1983): \textit{A Logit Model of Brand Choice Calibrated on Scanner Data}, Marketing Science, Vol. 2, No. 3, pp. 203-238.}

\subsubsection{1.1.3.3 — Non-Linear Relationships and Complex Interaction Effects}

Consumer behavior exhibits substantial non-linearity incompletely captured by traditional linear methodologies.\footnote{Hastie, T., Tibshirani, R., Friedman, J. (2009): \textit{The Elements of Statistical Learning}, 2nd Edition, Springer, New York, p. 237.} Small demographic differences can lead to giant behavioral gaps that univariate frameworks cannot properly model.

\subsubsection{1.1.3.4 — Absence of Optimality Guarantees}

Traditional segmentation imposes externally derived structures (like predetermined age ranges or RFM quintiles) rather than identifying empirically optimal partitions derived from actual patterns in behavioral data.

% ==============================================================================
% SECTION 1.2
% ==============================================================================
\section{1.2 — Theoretical Frameworks of Customer Value: Conceptual Framework and Probabilistic Models}

\subsection{1.2.1 — General Introduction to Customer Value Concepts}

Effective customer relationship management (CRM) constitutes a major strategic challenge for contemporary organizations. At the heart of this issue lies a fundamental question: how can one rigorously and predictively evaluate the value that a customer represents for the enterprise? This question has prompted practitioners and researchers to develop sophisticated conceptual and methodological frameworks for measuring what is commonly termed Customer Lifetime Value (CLV).

Traditionally, customer segmentation approaches relied on relatively simple criteria. However, the evolution of computational capabilities and the increasing availability of detailed transactional data have enabled the development of more granular and behavioral approaches. The concept of CLV has thus emerged as a central pillar of CRM strategy, enabling companies to identify and act upon the heterogeneity that characterizes their customer base.

\subsection{1.2.2 — Conceptual Definition and Economic Foundations of CLV}

Customer Lifetime Value represents the present value of expected future cash flows attributable to the customer-company relationship throughout its duration. Formally, CLV can be expressed as:

\begin{equation}
CLV = \sum_{t=0}^{T} \frac{(R_t - C_t)}{(1+d)^t}
\label{eq:clv_general}
\end{equation}

where:
\begin{itemize}
\item $R_t$ represents revenues generated by the customer in period $t$,
\item $C_t$ represents costs incurred to serve the customer in period $t$,
\item $d$ represents the discount rate reflecting the time value of money,
\item $T$ represents the expected duration of the customer relationship.
\end{itemize}

This formulation, grounded in microeconomic theory, provides a formal framework within which CLV can be estimated and operationalized. CLV serves as both a descriptive metric and a decision-making tool, enabling enterprises to concentrate acquisition and retention efforts on customer segments demonstrating the highest potential.

\subsection{1.2.3 — Probabilistic and Econometric Models for CLV Prediction}

The modeling of CLV has led to the development of several classes of models based on their theoretical foundation:

\subsubsection{1.2.3.1 — RFM Models and Descriptive Approaches}

The RFM (Recency, Frequency, Monetary) approach represents a heuristic method based on the hypothesis that these three variables constitute good predictors of future purchase behavior.\footnote{Hughes, A.M. (1996). "The Customer Loyalty Solution". New York: McGraw-Hill.} Although simple, it does not provide rigorous probabilistic estimates and assumes a linear relationship that may not be empirically verified.

\subsubsection{1.2.3.2 — Probabilistic Models in Contractual Contexts}

In contexts like subscriptions, probabilistic models based on Poisson processes and Markov chains have proven effective.\footnote{Soufiani, H., Jedidi, K., \& Degeratu, A.M. (2012). "A model of simultaneous diagonal and off-diagonal asymmetries in consumer preferences". Marketing Science, 31(6), 918-940.} The Exponential-Gamma (EG) model formalizes churn as a stochastic process governed by a Gamma-distributed rate parameter.

\subsubsection{1.2.3.3 — Non-contractual Probabilistic Models: The BG/NBD}

In non-contractual contexts where churn is not directly observable, the BG/NBD (Beta-Geometric/Negative Binomial Distribution) model is the standard.\footnote{Fader, P.S., Hardie, B.G.S., \& Lee, K.L. (2005). "Counting Your Customers the Easy Way". Journal of Marketing Research, 42(2), 139-155.} It rests on five assumptions:
\begin{enumerate}
\item Transaction numbers follow a Poisson distribution with rate $\lambda$.
\item Each customer has an unconditional probability $p$ of becoming inactive after each transaction.
\item Heterogeneity in transaction rate $\lambda$ follows a Gamma distribution.
\item Heterogeneity in churn probability $p$ follows a Beta distribution.
\item $\lambda$ and $p$ are independently distributed.
\end{enumerate}

\subsubsection{1.2.3.4 — Gamma-Gamma Models for Monetary Value}

Parallel to transaction prediction, the Gamma-Gamma model models the average monetary value per transaction.\footnote{Colombo, R., \& Morrison, D.E. (1989). "A Brand Switching Model with Implications for Marketing Expenditures". Marketing Science, 8(1), 89-99.} The joint distribution of expected monetary value $m_x$ and observed transactions $x$ is:
\begin{equation}
E[M|x, m_x] = \frac{(r+1)m_x}{r}
\label{eq:gamma_gamma}
\end{equation}

\subsection{1.2.4 — Relationship Between CLV, Customer Retention, and Profitability}

Several empirical studies have documented that the relationship between CLV and realized profitability depends crucially on cost structure.\footnote{Reinartz, W.J., Thomas, J.S., \& Kumar, V. (2005). "Balancing Acquisition and Retention Resources to Maximize Customer Profitability". Journal of Marketing, 69(1), 63-79.} Retention is justified only if it generates additional value offsetting the costs incurred. The relationship is often non-monotonic; concentrated effort on segments where retention effectively generates an increase in CLV is required.\footnote{Reinartz, W.J., \& Kumar, V. (2000). "On the profitability of long-life customers". Journal of Marketing, 64(4), 17-35.}

% ==============================================================================
% SECTION 1.3
% ==============================================================================
\section{1.3 — Sales Forecasting: Foundations and Methods}
\label{sec:1.3}

Sales forecasting designates the quantitative process by which an organisation estimates the volume and value of future sales over a defined time horizon. In the retail sector, it constitutes one of the most strategically significant analytical activities, since virtually every major operational decision depends on the quality of demand estimates.\footnote{G. Fildes, R. Kolassa, C. Ma: ``Retail Forecasting: Research and Practice'', \textbf{International Journal of Forecasting}, Vol. 38, No. 4, 2022, p. 1493.}

\subsection{1.3.1 — Importance of Forecasting in Retail Management}
\label{subsec:1.3.1}

Retailing is characterised by volatile consumer demand and short product life cycles. Under such conditions, accurate sales forecasting is a fundamental driver of operational efficiency.

\subsubsection{1.3.1.1 — The Strategic Role of Demand Forecasting}

Demand forecasts form the empirical backbone of investment decisions. Sales forecasting and effective demand planning positively affect the performance of the entire supply chain.\footnote{A. Abbasimehr, M. Shabani, M. Yousefi: ``A Comparative Study of Demand Forecasting Models'', \textbf{Operations Research Forum}, Vol. 3, No. 56, 2022, p. 2.} Forecasts must be produced simultaneously at multiple levels of granularity, giving rise to hierarchical forecasting frameworks.\footnote{R.J. Hyndman, G. Athanasopoulos: \textbf{Forecasting: Principles and Practice}, 2nd ed., 2018, p. 299.}

\subsubsection{1.3.1.2 — Operational Implications of Forecast Accuracy}

The consequences of poor focus propagate across the entire chain. Overestimation generates excess inventory, while underestimation produces stockouts and lost revenue.\footnote{S. Chopra, P. Meindl: \textbf{Supply Chain Management}, 6th ed., 2016, p. 187.} Accurate forecasting governs:
\begin{itemize}
  \item \textbf{Inventory management:} Calibration of safety stock.
  \item \textbf{Supply-chain coordination:} Reduction of the bullwhip effect.
  \item \textbf{Workforce scheduling:} Rosters centered on transaction volume.
  \item \textbf{Promotional planning:} Embedding uplift into forward-looking models.\footnote{G. Verstraete et al.: ``A Data-Driven Framework for Predicting Weather Impact'', Vol. 48, 2019, p. 170.}
\end{itemize}

\subsubsection{1.3.1.3 — Forecast Accuracy Metrics}

The quality of a model is universally assessed through standard error metrics. Table~1.1 defines the four most widely used indicators.

\begin{table}[h!]
\centering
{\bfseries Table 1.1: Standard Forecast Accuracy Metrics}\\[6pt]
\renewcommand{\arraystretch}{1.8}
\begin{tabular}{p{3.6cm} p{5.2cm} p{5.5cm}}
\toprule
\textbf{Metric} & \textbf{Formula} & \textbf{Interpretation} \\
\midrule
Mean Absolute Error (MAE)
  & $\displaystyle\frac{1}{n}\sum_{t=1}^{n}\lvert y_t - \hat{y}_t\rvert$
  & Average deviation in the same unit as sales. \\[6pt]
Root Mean Squared Error (RMSE)
  & $\displaystyle\sqrt{\frac{1}{n}\sum_{t=1}^{n}(y_t-\hat{y}_t)^{2}}$
  & Penalises large errors more heavily. \\[6pt]
Mean Absolute Percentage Error (MAPE)
  & $\displaystyle\frac{100}{n}\sum_{t=1}^{n} \left\lvert\frac{y_t-\hat{y}_t}{y_t}\right\rvert$
  & Scale-independent benchmarking. \\[6pt]
Symmetric MAPE (sMAPE)
  & $\displaystyle\frac{200}{n}\sum_{t=1}^{n} \frac{\lvert y_t-\hat{y}_t\rvert}{\lvert y_t\rvert+\lvert\hat{y}_t\rvert}$
  & Corrects the asymmetry of MAPE. \\
\bottomrule
\end{tabular}\\[5pt]
{\small\textit{Source: Compiled by the author, based on Hyndman \& Athanasopoulos (2018), p. 53.}}
\end{table}

\subsection{1.3.2 — Classical Methods: Time Series, Regression, and Exponential Smoothing}
\label{subsec:1.3.2}

Retailers relied on well-established toolkit of statistical methods including Moving Averages (MA):
\begin{equation}
  \hat{y}_{t+1} \;=\; \frac{1}{k}\sum_{i=0}^{k-1} y_{t-i}
  \label{eq:ma}
\end{equation}

Exponential smoothing addresses the limitations of MA by assigning geometrically decreasing weights. 
\begin{itemize}
    \item \textbf{Simple Exponential Smoothing (SES):} $\hat{y}_{t+1} = \alpha\,y_t + (1-\alpha)\,\hat{y}_t$.
    \item \textbf{Holt's Double Smoothing:} Extends to linear trends.
    \item \textbf{Holt-Winters Triple Smoothing:} Incorporates seasonality components.
\end{itemize}

ARIMA $(p, d, q)$ specification combines Autoregression, Integration (differencing), and Moving Average error. The seasonal extension (SARIMA) captures periodic demand cycles.\footnote{P. Ramos et al., Vol. 34, 2015, p. 152.}

Table 1.2 consolidates these methods:

\begin{table}[h!]
\centering
{\bfseries Table 1.2: Comparative Overview of Classical Sales Forecasting Methods}\\[6pt]
\renewcommand{\arraystretch}{1.65}
\resizebox{\textwidth}{!}{%
\begin{tabular}{p{2.8cm} p{3.2cm} p{1.9cm} p{2.3cm} p{2.3cm} p{1.8cm} p{3cm}}
\toprule
\textbf{Method} & \textbf{Mechanism} & \textbf{Trend} & \textbf{Season.} & \textbf{Exogen.} & \textbf{Data} & \textbf{App.} \\
\midrule
SMA & Arithmetic Mean & No & No & No & Low & Stable demand \\
SES & Decay-weighted & No & No & No & Low & Level-stationary \\
Holt & Double Smoothing & Yes & No & No & Low & Trending demand \\
Holt-Winters & Triple ES & Yes & Yes & No & Medium & Seasonal retail \\
ARIMA & AR+I+MA & Yes & No & No & Medium & Time-series \\
SARIMA & Seasonal ARIMA & Yes & Yes & No & Medium & Monthly/Weekly \\
Regression & Causal model & Yes & Yes & Yes & Medium & Promo effects \\
Prophet & Decomposable & Yes & Yes & Yes & Low & Retail KPIs \\
\bottomrule
\end{tabular}}
\end{table}

\subsection{1.3.3 — Contribution of Machine Learning to Sales Forecasting}
\label{subsec:1.3.3}

Machine-learning methods have emerged as the dominant response to retail complexity.
Ensemble tree-based methods aggregate predictions:
\begin{itemize}
    \item \textbf{Random Forest (RF):} constructs $B$ trees on independent bootstrap samples. Forecast is $\hat{y} = \frac{1}{B}\sum f_b(\mathbf{x})$.
    \item \textbf{XGBoost:} trains trees sequentially correcting residual errors.
\end{itemize}

\textbf{Deep Learning Architectures}:
\begin{itemize}
    \item \textbf{ANN (MLP):} approximate arbitrary continuous functions.
    \item \textbf{LSTM Networks:} incorporate three gating mechanisms (forget, input, output) to retain long-range information.
    \item \textbf{CNN:} effective at detecting short-cycle spikes from temporal windows.
\end{itemize}

Hybrid models like CNN-LSTM or ATES (Adaptive Temporal ES) consistently demonstrate the highest accuracy on retail benchmarks.\footnote{M.I. Efat et al., Vol. 158, 2024, p. 12.}

% ==============================================================================
% CHAPTER CONCLUSIONS (CONSOLIDATED)
% ==============================================================================
\section*{Conclusion of Chapter I}
\addcontentsline{toc}{section}{Conclusion of Chapter I}

This chapter has established the foundational conceptual and methodological framework for understanding customer segmentation, CLV, and forecasting as strategic disciplines within data-driven retail.

We have demonstrated that customer behavior and sales dynamics are not independent: purchasing patterns identified through RFM segmentation constitute predictive signals that enrich demand forecasting models. The probation and econometric models presented offer sophisticated tools for estimating CLV in varied empirical contexts, with the BG/NBD model emerging as a status standard for non-contractual retail.

Furthermore, machine-learning methods have bridged the gap where classical statistical models fail to capture the complex, non-linear interactions characterizing modern retail demand. These foundations collectively motivate the methodological framework of Chapter 2, which operationalizes the RFM-ML segmentation approach and selects the forecasting architecture best suited to the data and objectives of this study.

% ==============================================================================
% BIBLIOGRAPHY (CONSOLIDATED)
% ==============================================================================
\clearpage
\section*{Bibliography}
\addcontentsline{toc}{section}{Bibliography}

\subsection*{I. Books}
\begin{itemize}
    \item Bellman, R. E. (1961): \textit{Adaptive Control Processes: A Guided Tour}, Princeton University Press, New Jersey.
    \item Blattberg, R. C., Getz, G., \& Thomas, J. S. (2001): \textit{Customer Equity}, Harvard Business School Press.
    \item Breiman, L.: \textbf{Random Forests — Methods and Applications}, Kluwer, 2001, 128 p.
    \item Chopra, S., Meindl, P.: \textbf{Supply Chain Management}, 6th ed., Pearson, 2016.
    \item Hughes, A. M. (2005): \textit{Strategic Database Marketing}, 3rd Edition, AMA, Chicago.
    \item Hughes, A.M. (1996). \textit{The Customer Loyalty Solution}. New York: McGraw-Hill.
    \item Hyndman, R.J., Athanasopoulos, G.: \textbf{Forecasting: Principles and Practice}, 2nd ed., 2018.
    \item Kotler, P., Keller, K. L. (2016): \textit{Marketing Management}, 15th Edition, Pearson.
    \item Schiffman, L. G., Kanuk, L. L. (2010): \textit{Consumer Behavior}, 10th Edition, Prentice Hall.
    \item Stevenson, W.J.: \textbf{Operations Management}, 12th ed., McGraw-Hill, 2014.
    \item Wedel, M., Kamakura, W. A. (2000): \textit{Market Segmentation}, 2nd Edition, Kluwer.
\end{itemize}

\subsection*{II. Journal Articles and Proceedings}
\begin{itemize}
    \item Abbasimehr, A. et al.: ``A Comparative Study of Demand Forecasting Models'', \textbf{Operations Research Forum}, 2022.
    \item Ahmadov, T., Helo, P.: ``Applying Deep Neural Networks'', \textbf{Expert Systems with Applications}, 2023.
    \item Alon, I. et al.: ``Forecasting Aggregate Retail Sales'', \textbf{JRCS}, 2001.
    \item Aras, S. et al.: ``Comparative Study on Retail Sales Forecasting'', \textbf{JBEM}, 2017.
    \item Colombo, R., \& Morrison, D.E. (1989). "A Brand Switching Model". \textit{Marketing Science}, 8(1).
    \item DeCarlo, T. E. (2005): \textit{Suspicion of Ulterior Motives}, JCP, Vol. 15.
    \item Efat, M.I. et al.: ``Adaptive Temporal Exponential Smoothing'', \textbf{Applied Soft Computing}, 2024.
    \item Fader, P.S., Hardie, B.G.S., \& Lee, K.L. (2005). "Counting Your Customers". \textit{JMR}, 42(2).
    \item Fildes, G. et al.: ``Retail Forecasting'', \textbf{IJF}, 2022.
    \item Haque, K. et al.: ``Sales Forecasting Using Hybrid Neural Networks'', \textbf{PLOS ONE}, 2025.
    \item Ramos, P. et al.: ``Performance of State Space and ARIMA Models'', \textbf{RCIM}, 2015.
    \item Reinartz, W.J. et al. (2005). "Balancing Acquisition and Retention". \textit{Journal of Marketing}, 69(1).
    \item Smyl, S.: ``A Hybrid Exponential Smoothing and RNN for Time Series'', \textbf{IJF}, 2020.
    \item Wind, Y. (1978): \textit{Issues and Advances in Segmentation Research}, JMR, Vol. 15.
\end{itemize}

\subsection*{III. Theses and Official Reports}
\begin{itemize}
    \item Almesfer, A.: ``Forecast-Driven Inventory Management'', Master Thesis, MIT, 2023.
    \item WJARR: ``Optimizing Inventory Management and Demand Forecasting'', Vol. 20, 2023.
\end{itemize}

\end{document}
